\subsection{Biểu đồ Use Case tổng quát}

Trên cơ sở các tác nhân đã xác định ở mục trước, toàn bộ chức năng nghiệp vụ của hệ thống được mô hình hóa dưới dạng \textit{Use Case Diagram} tổng quát. Biểu đồ này thể hiện phạm vi tương tác giữa từng tác nhân và các nhóm chức năng tương ứng với phân quyền, đồng thời làm rõ ranh giới giữa các chức năng yêu cầu xác thực và các chức năng công khai phục vụ kiểm chứng minh bạch.

\begin{figure}[!ht]
    \centering
    \includegraphics[width=0.58\textwidth, keepaspectratio]{Diagrams/Usecase/TongQuanHeThong.pdf}
    \caption{Biểu đồ Use Case tổng quát của hệ thống}
    \label{fig:uc-tong-quat}
\end{figure}

Các chức năng được phân chia theo năm nhóm chính, mỗi nhóm tương ứng với một phạm vi nghiệp vụ và mức phân quyền riêng biệt:

\begin{itemize}
    \item \textbf{Nhóm chức năng dùng chung:} Bao gồm các thao tác xác thực phiên làm việc và xem thông tin tài khoản cá nhân, áp dụng cho cả quản trị viên và cử tri sau khi đăng nhập hợp lệ.

    \item \textbf{Nhóm chức năng quản lý người dùng:} Dành riêng cho quản trị viên, bao gồm tạo tài khoản đơn lẻ hoặc hàng loạt, tìm kiếm, xem chi tiết, cập nhật thông tin và quản lý trạng thái kích hoạt của tài khoản cử tri.

    \item \textbf{Nhóm chức năng quản lý cuộc bầu cử:} Cũng thuộc quyền của quản trị viên, bao gồm tạo cuộc bầu cử, quản lý danh sách ứng viên và cử tri tham gia, điều hành cuộc bầu cử (bắt đầu, đóng) và theo dõi danh sách phiếu đã nộp.

    \item \textbf{Nhóm chức năng bỏ phiếu của cử tri:} Cho phép cử tri xem danh sách các cuộc bầu cử mà mình được phép tham gia, thực hiện bỏ phiếu thông qua cơ chế chữ ký mù và tiết lộ phiếu sau khi cuộc bầu cử kết thúc.

    \item \textbf{Nhóm chức năng công khai:} Phục vụ tác nhân khách ẩn danh, không yêu cầu xác thực, bao gồm xem danh sách cuộc bầu cử công khai, xác minh phiếu bầu, xem kết quả kiểm phiếu và kiểm toán toàn bộ số phiếu trên sổ cái blockchain.
\end{itemize}
